\documentclass[10pt]{article}
\usepackage{fontspec}
\setmainfont{Noto Serif CJK SC}
\setsansfont{Noto Sans CJK SC}
\XeTeXlinebreaklocale "zh"
\XeTeXlinebreakskip=0pt plus 1pt
\usepackage[a4paper,margin=20mm]{geometry}
\usepackage{amsmath,amssymb,booktabs,microtype,url,xcolor,float}
\usepackage[section]{placeins}
\usepackage[hidelinks]{hyperref}
\usepackage[numbers,sort&compress]{natbib}
\setlength{\emergencystretch}{2em}
\renewcommand{\abstractname}{摘要}
\renewcommand{\refname}{参考文献}
\renewcommand{\tablename}{表}
\newcommand{\DevelopmentIncidentCount}{403}
\newcommand{\DevelopmentCausalCount}{256}
\newcommand{\DevelopmentMethodRate}{100.0}
\newcommand{\DevelopmentComparatorRate}{51.6}
\newcommand{\DevelopmentGain}{48.4}
\newcommand{\DevelopmentControlRate}{100.0}
\newcommand{\DevelopmentRollbackRate}{100.0}
\newcommand{\ConfirmationRosterCount}{74}
\newcommand{\SuccessfulSourceCount}{71}
\newcommand{\ValidCausalCount}{226}
\newcommand{\ValidControlCount}{71}
\newcommand{\PrimaryMethodRate}{100.0}
\newcommand{\PrimaryComparatorRate}{49.6}
\newcommand{\PrimaryGain}{50.4}
\newcommand{\PrimaryCILow}{46.6}
\newcommand{\PrimaryCIHigh}{54.6}
\newcommand{\PrimaryP}{0.00001}
\newcommand{\ControlLoss}{0.0}
\newcommand{\ConfirmationStatus}{confirmation gate passed}
\newcommand{\MeanAdapterActions}{7.14}
\newcommand{\MaxAdapterActions}{110}
\newcommand{\MeanRelyAdapterActions}{18.89}


\title{修复边界而非技能：面向 LLM Agent 技能组合的契约引导适配器合成}
\author{匿名作者}
\date{匿名稿，2026 年 9 月}

\begin{document}
\maketitle

\begin{abstract}
可复用 Skill 使 LLM Agent 能够保留成功程序，但各自具有成功见证的组件仍可能因
组合边界上的依赖条件、角色绑定、活跃状态或验证不一致而失败。修改组件会扩大回归
范围，回滚整个链则会丢弃无关变更。本文研究一个更窄的选择：保持组件产物不变，只在
已诊断边界编译一个可执行中介。我们提出契约引导边界适配器合成（Contract-Guided
Boundary Adapter Synthesis，CGBAS），结合类型化 rely/guarantee 契约、下游
活跃性和保留的成功来源证据，插入缺失的见证前缀、特化消费端角色，或在执行前守卫
破坏性调用。一次失败的早期机制修订被完整保留。在 \DevelopmentIncidentCount{} 个
公开开发事件上完成修订后，我们先冻结方法、终点、最强对照、三个模型版本和分析器，
再选择确认任务。独立确认使用 \ConfirmationRosterCount{} 个与既往任务标识零重叠的
ALFWorld 和 ScienceWorld 任务，形成 \ValidCausalCount{} 个状态匹配因果冲突和
\ValidControlCount{} 个健康对照。CGBAS 的约束原生修复成功率为
\PrimaryMethodRate\%，冻结最强局部对照为 \PrimaryComparatorRate\%，配对提升
\PrimaryGain{} 个百分点（95\% 按任务族/环境分层、按任务聚类的 bootstrap 区间
\PrimaryCILow--\PrimaryCIHigh；双侧任务符号随机化检验 $p=\PrimaryP$），健康对照
损失为 \ControlLoss{} 个百分点。结论有明确边界：CGBAS 假设诊断正确且成功来源
证据仍然有效；一个适配器调用也可能包含很长的见证前缀。本文不声称能修复任意生产
故障，也不证明局部适配总是优于重新规划。
\end{abstract}

\noindent\textbf{关键词：}LLM Agent；Agent Skill；技能组合；接口修复；契约式设计；
来源追踪；ALFWorld；ScienceWorld

\section{引言}

LLM Agent Skill 将可复用程序、脚本、工具约定和验证说明封装为软件产物。复用使一次
成功轨迹可以在后续任务中被选择与组合，但组合也会产生孤立评测无法发现的故障：消费
调用依赖的状态可能未由前驱建立；两个调用可能把同一语义角色绑定到不同对象；中间
调用可能破坏后续仍需要的状态；验证结论也可能在某个效果发生后失效。

已有研究将 Skill 子集、数量和顺序联合建模\citep{zhao2026skillcomposer}，把 Skill
编译为带类型且具有局部修复操作的执行图\citep{xia2026grasp}，或通过类型化契约诊断
技能库兼容性债务\citep{pu2026skillops}。来源审计把复用行为连接到操作证据
\citep{chen2026multitrace}，持久决策历史则支持跨会话演化
\citep{li2026skillhone}。这些工作引出一个容易与规划混淆的维护问题：当具体接口已经
被诊断后，维护者是否必须改写组件或替换整个计划，还是可以直接修复接口？

这一差别具有实际意义。修改整个 Skill 会使该 Skill 所有既往用途的证据失效；回滚
整条链虽然恢复见证版本，却会丢弃未受影响的独立变更；拒绝执行能避免损害，但会损失
能力。边界适配器提供第四种选择：保留组件身份，只调解已诊断交接。但只有当它在原生
环境中成功、动作有依据、来源可追踪且不损害健康链时，这种选择才成立。

CGBAS 的输入是有序候选链、最小冲突诊断、公开任务目标和来源关联的成功见证链。它
最多执行一次边界操作：重放固定消费端之前被遗漏的见证前缀，用角色一致的见证消费
调用包裹错误调用，或在非目标破坏调用发生前抑制它。健康链保持不变。因此，适配器是
对接口的可执行补丁，而不是重新生成的整条计划。

评测针对三类有效性风险设计。第一，静态契约闭合不能替代原生任务成功。第二，LLM
对照必须公平：三个固定版本的开放指令模型接收相同事件、诊断、见证目录和四类允许
操作。第三，开发决策与确认必须隔离：既往诊断论文的公开记录只用于两次机制修订，
失败过程全部保留；随后对方法、对照、模型、终点、门槛、环境执行器和分析代码进行
哈希锁定，再通过与内容无关的方式选择零重叠任务。

本文贡献如下：
\begin{itemize}
  \item 提出具有明确不变量的边界修复问题：保持所有未受影响调用身份不变，并最多
  引入一个适配器调用；
  \item 提出 CGBAS，把 rely/guarantee、活跃性、角色、验证和成功来源证据编译为
  三类可执行修复操作；
  \item 在原生环境中与无修复、无来源/全局活跃性的成对契约启发式、整链回滚和三个
  输出模式匹配的局部 LLM 对照进行因果比较；
  \item 在 ALFWorld 与 ScienceWorld 上采用可审计的开发/确认隔离，报告无效源任务、
  失败修订、任务聚类推断、动作成本和单张 RTX 3090 复现材料。
\end{itemize}

研究问题为：\textbf{RQ1}，契约与来源引导的边界适配器能否在独立选择的组合冲突上
恢复原生成功？\textbf{RQ2}，它能否超过冻结的最强可部署局部对照且不损害健康链？
\textbf{RQ3}，哪些假设与动作成本限制了结论？

\begin{figure}[t]
\centering
\fbox{\begin{minipage}{0.94\linewidth}\centering
公开开发 $\rightarrow$ 两次机制修订 $\rightarrow$ 方法/对照/分析哈希锁\\[3pt]
$\Downarrow$\\[3pt]
内容无关的零重叠名单 $\rightarrow$ 成功源见证 $\rightarrow$ 冻结确定性与 LLM 程序\\[3pt]
$\Downarrow$\\[3pt]
状态匹配候选/修复执行 $\rightarrow$ 七臂原生重放 $\rightarrow$ 任务聚类推断
\end{minipage}}
\caption{证据生成顺序。确认名单、源事件和全部待评程序冻结前，不产生确认结果。}
\label{fig:evidence-order}
\end{figure}

\section{相关工作}

\paragraph{Skill 选择与组合。}
Generative Skill Composition 将选择、数量和顺序定义为一个结构化预测问题
\citep{zhao2026skillcomposer}。GraSP 在检索与执行之间加入编译层，用带类型 DAG
表示依赖并提供节点验证与局部修复\citep{xia2026grasp}。SkillFuzz 用抽取的契约在
执行前优先搜索 Skill 交互\citep{hu2026skillfuzz}。本文研究链条中更晚且更窄的
环节：选择和顺序固定，具体边界已经被诊断，终点是单个中介能否在状态匹配的原生任务
中恢复成功而不改写组件。

\paragraph{生命周期维护与审计。}
SkillOps 将持续积累的兼容性与验证缺陷视为 Skill 技术债务
\citep{pu2026skillops}。SkillHone 跨会话保留诊断、修订、证据和被拒绝方案
\citep{li2026skillhone}。SkillAudit 通过配对轨迹和固定结构验证器，在演化过程中不
读取隐藏奖励也能指导修复\citep{gao2026skillaudit}。SkillCoach 则区分过程质量与
结果成功\citep{zhu2026skillcoach}。本文把这些原则落实为保留决策历史、配对执行，
并同时报告局部性、动作支持、来源和成本；不同点是修复执行边界而非 Skill 文档。

\paragraph{契约与来源。}
前置/后置条件使程序假设显式化\citep{hoare1969axiomatic,meyer1992contract}；
rely/guarantee 推理区分组件依赖与组件在受干扰状态中建立的事实
\citep{jones1983rely}。本文只把这些思想作为文本模拟器调用的操作抽象，不声称验证
任意并发程序。SkillTrace 表明 Skill 来源分布在表达、实现和操作轨迹中
\citep{chen2026multitrace}。CGBAS 更窄地使用操作轨迹：新增命令或角色包裹必须引用
成功见证，除非原生环境本身接受该命令。

\section{问题定义}

有序 Skill 组合为 $C=(c_1,\ldots,c_n)$。每个调用包含不可变动作 $A_i$、标识、来源
见证和契约
\begin{equation}
K_i=(R_i,G_i,D_i,B_i,V_i),
\end{equation}
其中 $R_i$ 是依赖集合，$G_i$ 是保证谓词，$D_i$ 是被破坏谓词，$B_i$ 是语义角色
绑定，$V_i$ 是已验证谓词。若 $R_i\subseteq S_{i-1}$ 且绑定兼容，抽象状态按
\begin{equation}
S_i=(S_{i-1}\setminus D_i)\cup G_i
\end{equation}
更新；逆向活跃性分析标记任一后继或目标完成仍需使用的谓词。

本文覆盖四类诊断冲突：消费端需要不存在谓词的\emph{依赖缺口}；同一角色被重新绑定
到不同对象的\emph{角色错配}；适用效果破坏活跃谓词的\emph{状态破坏}；以及该破坏
发生在谓词检查后的\emph{陈旧验证}。不存在这些冲突且建立目标完成谓词的链为健康链。

令 $W=(w_1,\ldots,w_m)$ 为同一任务的成功来源链。合法变换 $T(C,W)$ 只能在已诊断
边界新增至多一个适配器，未被包裹的调用保持标识和内容不变。主终点为约束原生修复成功：
\begin{equation}
Y=\mathbb{1}[\text{原生目标成功}\land\text{局部性}\land
\text{来源完整}\land\text{适配器命令有支持}].
\end{equation}
适配器命令只要精确来自保留的成功证据或被原生环境接受即视为有支持。组件中原本存在
但被环境拒绝的探索命令单独报告，不反向归因给适配器。

\section{契约引导边界适配器}

\subsection{诊断与修复槽位}
CGBAS 对候选链进行一次扫描，维护抽象状态、活跃绑定和验证位置。首个冲突返回类别、
证据和最小调用对。诊断由前序工作提供而非本文贡献；本文干预从该槽位开始。

\subsection{缺失生产前缀}
对于消费调用 $c_j$ 之前的依赖缺口，CGBAS 在 $W$ 中定位同一不可变消费调用标识，
选择其之前且候选前缀尚未执行的有序见证调用，把动作连接为一个适配器，并合并契约与
见证摘要。它不只选择满足一个谓词的最短调用，因此能保留清洁、冷却、加热或测量等
目标相关处理。由此，一个调用级修改也可能包含很多环境动作。

\subsection{见证角色包裹}
对于角色错配，CGBAS 将活跃生产端绑定与成功见证中的目标完成调用比较，选择非目的地
角色与活跃绑定一致的最短见证消费调用，在错误槽位用它进行包裹，并记录来源与被替换
标识。组件产物仍然保留，只有本次调用被调解。

\subsection{执行前保持守卫}
若破坏性动作会终止 episode 或不可逆改变世界，事后补偿将不可能。对于状态破坏和
陈旧验证，CGBAS 在执行前检查效果适用性、是否满足目标以及下游活跃性。被诊断为非
目标的破坏调用由零动作守卫包裹；其原标识、动作、效果证据和见证保留在元数据中。

\subsection{复杂度与不变量}
诊断相对于调用数和契约集合运算为线性过程；见证前缀和角色选择只扫描见证链；守卫只
需一次定位。方法不调用 LLM，最多改变一个调用边界，健康链保持不变。这是调用级局部
性，而不是常数动作成本保证。

\section{实验方法}

\subsection{开发与修订历史}
开发使用既往诊断研究公开的 403 个事件，包括 256 个因果冲突和 147 个健康对照，
覆盖 ALFWorld 与 ScienceWorld\citep{shridhar2021alfworld,wang2022scienceworld}。
修订 1 插入最短谓词生产调用，并在破坏后补偿；其冲突成功率为 70.6\%，低于冻结的
75\% 门槛。失败表明它遗漏了目标相关前缀，且 ScienceWorld 破坏调用会提前终止。
修订 2 将破坏干预前移，并保留完整缺失见证前缀。第一版代码、输出、399 条原生轨迹
以及四个被遗漏的弃权均保存在产物中。

确认前，修订 2 与无修复、无来源和全局活跃性的成对契约启发式、整链回滚，以及
Qwen3-4B-Instruct-2507、Phi-4-mini-instruct、Mistral-7B-Instruct-v0.3 三个
局部模型比较。模型接收完全相同的公开输入，输出只允许 no-op、有序见证前缀、见证
消费替换、保持守卫或弃权；无效输出和弃权执行原候选并记为失败。最初沿用修订 1 的
LLM 输出模式被识别为不公平，相关完整/部分结果被保留但排除；三个模型随后都在与
修订 2 匹配的模式下重跑。

按开发约束成功率自动选择最强可部署局部方法为主对照。整链回滚作为普通成功上界报告，
但因违反局部性不参与主对照选择。开发门槛包括至少 180 个因果事件和 60 个对照、
至少 75\% 约束成功、相对最强对照至少 15 点提升、每类和两个环境均有正向修复收益、
健康对照损失不超过 2 点，以及至少 95\% 适配器命令支持率。

\subsection{独立确认}
选择前先合并论文 1--12 与本文开发所接触的任务标识。salted hash 仅使用环境、任务族、
标识、官方划分、兼容标记和排除状态排序。冻结名单包含五个家族各 8 个未触及
ALFWorld train 实例，以及全部 34 个合格且未触及的 ScienceWorld use-thermometer
test 变体。ALFWorld 部分支持任务级零重叠迁移结论，不代表官方测试集表现。

名单冻结后才采集原生专家/金标准轨迹。失败源任务保留并排除，不替换。每个成功源
产生一个健康对照和四类可能的单边界变异：遗漏生产前缀；使用同族但不同 subject 的
donor 消费调用；插入适用的非目标破坏调用；或在破坏前加入验证。变异规格和公开事件
在任何候选/修复执行前冻结。

随后冻结所有确定性与 LLM 程序，才从相同原生初始状态执行候选和匹配修复。因果标签
要求源证据成功、候选失败、匹配修复成功、两次初始状态匹配且无技术错误；健康对照
要求源与未修改候选都成功。无效变异直接排除，不重标为健康对照。各程序臂在全部计划
事件上独立重放，技术错误按失败保留。
图~\ref{fig:evidence-order} 概括这一顺序，各阶段文件哈希使边界可审计。

\subsection{统计推断与门槛}
主比较为 CGBAS 与冻结最强对照在有效因果冲突上的配对约束成功差。一个源任务可产生
多个冲突，因此推断按任务聚类。我们使用 10,000 次按任务族/环境分层的任务 bootstrap
构造 95\% 区间，并执行 100,000 次双侧配对任务符号随机化检验
\citep{efron1994bootstrap}。确认至少需要 60 个源任务、180 个因果冲突、60 个对照和
每类 30 个样本；CGBAS 至少 75\%；相对最强对照至少提升 15 点、区间下界大于 0 且
$p\leq0.01$；相对无修复在每类和两个环境均为正；对照损失不超过 2 点；并通过局部性、
来源、初始状态、命令支持和技术错误审计。

\section{结果}

\subsection{开发门槛}
\begin{table}[t]
\centering
\small
\caption{公开开发结果。独立确认前按预声明规则冻结最强合格局部对照。}
\begin{tabular}{lrrr}
\toprule
方法 & 原生 (\%) & 约束 (\%) & 健康对照 (\%) \\
\midrule
CGBAS & 100.0 & 100.0 & 100.0 \\
Mistral-7B & 51.6 & 51.6 & 100.0 \\
Phi-4-mini & 46.5 & 46.5 & 100.0 \\
Qwen3-4B & 35.2 & 35.2 & 100.0 \\
Pairwise contract & 10.2 & 9.8 & 100.0 \\
No repair & 0.0 & 0.0 & 100.0 \\
Whole-chain rollback & 100.0 & 0.0 & 100.0 \\
\bottomrule
\end{tabular}
\end{table}


修订 2 在 \DevelopmentCausalCount{} 个因果冲突上的约束成功率为
\DevelopmentMethodRate\%。Mistral-7B 由规则自动选择为最强合格对照，成功率为
\DevelopmentComparatorRate\%，开发差值为 \DevelopmentGain{} 点。健康对照成功率
保持 \DevelopmentControlRate\%。整链回滚的普通成功率为
\DevelopmentRollbackRate\%，但不满足局部终点。

\subsection{独立主结果}
\begin{table}[t]
\centering
\small
\caption{独立确认结果。无效变异通过各方法共用的状态匹配标签规则排除。}
\begin{tabular}{lrrrr}
\toprule
方法 & $n$ & 原生 (\%) & 约束 (\%) & 健康对照 (\%) \\
\midrule
CGBAS & 226 & 100.0 & 100.0 & 100.0 \\
Mistral-7B & 226 & 49.6 & 49.6 & 100.0 \\
Phi-4-mini & 226 & 42.9 & 42.9 & 100.0 \\
Qwen3-4B & 226 & 36.7 & 36.7 & 100.0 \\
Pairwise contract & 226 & 8.4 & 7.5 & 100.0 \\
No repair & 226 & 0.0 & 0.0 & 100.0 \\
Whole-chain rollback & 226 & 100.0 & 0.0 & 100.0 \\
\bottomrule
\end{tabular}
\end{table}


在 \ConfirmationRosterCount{} 个名单任务中，\SuccessfulSourceCount{} 个产生有效源
轨迹；状态匹配规则保留 \ValidCausalCount{} 个因果冲突和 \ValidControlCount{} 个
对照。CGBAS 约束成功率为 \PrimaryMethodRate\%，冻结对照为
\PrimaryComparatorRate\%，配对差为 \PrimaryGain{} 点，95\% 聚类区间为
\PrimaryCILow--\PrimaryCIHigh，任务符号随机化检验 $p=\PrimaryP$。预声明确认状态为
\ConfirmationStatus{}。

\subsection{冲突类别、环境与成本}
\begin{table}[t]
\centering
\small
\caption{各冲突类别的约束原生成功率。}
\begin{tabular}{lrrrr}
\toprule
冲突类别 & $n$ & CGBAS (\%) & 冻结对照 (\%) & 无修复 (\%) \\
\midrule
依赖缺口 & 71 & 100.0 & 43.7 & 0.0 \\
角色错配 & 71 & 100.0 & 0.0 & 0.0 \\
状态破坏 & 42 & 100.0 & 100.0 & 0.0 \\
陈旧验证 & 42 & 100.0 & 92.9 & 0.0 \\
\bottomrule
\end{tabular}
\end{table}

\begin{table}[t]
\centering
\small
\caption{按原生环境划分的因果约束成功率。}
\begin{tabular}{lrrrr}
\toprule
环境 & $n$ & CGBAS (\%) & 冻结对照 (\%) & 无修复 (\%) \\
\midrule
ALFWorld & 128 & 100.0 & 64.1 & 0.0 \\
ScienceWorld & 98 & 100.0 & 30.6 & 0.0 \\
\bottomrule
\end{tabular}
\end{table}


分类型收益相对于无修复报告，整体 15 点门槛则始终相对于冻结最强对照。该区分避免
一个机械上较容易的类别出现对照平局时被总体平均掩盖。CGBAS 每个因果修复平均提出
\MeanAdapterActions{} 个适配器动作，最多 \MaxAdapterActions{} 个；依赖缺口平均为
\MeanRelyAdapterActions{} 个，说明“一个边界调用”和“很短的执行”并不等价。

\subsection{完整性审计}
\begin{table}[t]
\centering
\small
\caption{预声明确认门槛与完整性检查。}
\begin{tabular}{lr}
\toprule
冻结检查项 & 结果 \\
\midrule
minimum tasks & 通过 \\
minimum causal conflicts & 通过 \\
minimum controls & 通过 \\
minimum per conflict class & 通过 \\
minimum constrained success & 通过 \\
minimum gain over frozen comparator & 通过 \\
bootstrap interval above zero & 通过 \\
paired randomization p at most alpha & 通过 \\
positive repair gain every conflict class & 通过 \\
positive repair gain both environments & 通过 \\
control loss at most two points & 通过 \\
minimum 95 percent adapter command support & 通过 \\
all primary repairs local & 通过 \\
all primary provenance complete & 通过 \\
all primary initial states match & 通过 \\
no primary technical errors & 通过 \\
zero overlap roster & 通过 \\
\bottomrule
\end{tabular}
\end{table}


确认名单与既往接触任务交集为 0。所有程序产物先于候选和适配器结果生成；模型版本标识、
提示词、方法、对照、执行器、终点、门槛和分析哈希均与确认前锁一致。无效源与变异仍
保存在标签表，各臂的全部计划结果在共同有效性过滤前均保留。

\section{讨论}

结果的实际意义是有条件的：当维护者已经获得可靠的边界诊断和仍有效的成功来源证据
时，修复不必演化为整个 Skill 的修改。依赖缺口要恢复完整因果前缀，而不只是满足一个
抽象谓词；可能无法补偿的破坏必须在发生前守卫；角色故障可以用见证消费调用修复，而
无需自由生成命令。这些机制解释了即使得到同一诊断槽位，成对契约启发式与局部 LLM
选择仍可能不足。

整链回滚仍然重要。如果来源重放便宜、不需要保留独立变更，或诊断不确定，回滚可能
更合适。CGBAS 位于另一工程取舍点：它保留未受影响调用身份，并显式记录插入或抑制
内容。但调用局部性不能等同于运行经济性，因为缺失生产端可能需要很长的见证前缀。

该结果也支持一种生命周期架构：组合诊断器输出带证据的类型化边界；来源库提供兼容
见证切片；适配器编译器产生版本化调用包裹；发布门禁同时重放冲突和健康对照。不可变
组件与显式适配器比静默修改共享 Skill 更容易回滚、审计和后续清理。

\section{有效性威胁}

\paragraph{构造有效性。}
四类变异是从成功轨迹产生的受控接口回归，不代表生产环境中的发生率。因果标签要求
候选原生失败和匹配修复成功，但开发与确认使用相同变异操作，因此证据支持新任务迁移，
不支持发现新故障类型。

\paragraph{信息优势。}
CGBAS 获得正确诊断和成功见证目录，这适合保留来源证据的维护场景，却强于任意故障
后的可用信息。LLM 对照接收同一目录、诊断和操作集合，但它们也不是一般自主 Agent。

\paragraph{外部有效性。}
ALFWorld 和 ScienceWorld 提供确定性文本接口、原生目标和可重置状态。生产 API 可能
漂移、有副作用或缺少精确重置与结果 oracle。严格零重叠后，ScienceWorld 确认只剩
一个合格任务族；扩展到工具、网页和异步长链后才能提出更广结论。

\paragraph{统计有效性。}
同一源任务的派生冲突并不独立，因此全部推断按任务聚类并按任务族/环境分层。但 74 个
名单任务仍然有限，各类估计比主聚合结果更不精确。门槛和分析在确认前冻结，失败源未被
替换。

\paragraph{实现有效性。}
契约由成功轨迹编译，原生命令支持只对适配器新增动作审计。原 Skill 可能含有未出现在
环境动作列表中的探索命令；本文报告但不把它归因于适配器。哈希锁、追加式轨迹、初始
状态签名和独立复现包降低但不能消除实现风险。

\section{结论}

CGBAS 通过编译一个有来源证据的边界适配器，修复已经诊断的 LLM Agent Skill 接口。
它恢复遗漏的见证前缀、特化消费端角色，或在执行前阻止活跃状态被破坏，同时保留未受
影响组件身份。冻结的零重叠原生评测将该机制与无修复、只看成对契约的局部修复、局部性
约束下的整链回滚和三个输出模式匹配的 LLM 选择器比较。本文结论只适用于成功来源仍
有效的受控接口回归；在这一边界内，修复接口是改写可复用 Skill 之外的一种可审计选择。

\bibliographystyle{plainnat}
\bibliography{../references}
\end{document}
